\section{Introduction}

Flows in field space are an interesting tool that may
allow new insights to be gained into the physical
mechanisms described
by highly non-linear quantum field theories such as QCD.
The flow $B_{\mu}(t,x)$ of $\SUthree$ gauge fields
studied in this paper is
defined by the equations
\begin{align}
  \dot{B}_{\mu} &= D_{\nu}G_{\nu\mu},
  &
  \left.B_{\mu}\right|_{t=0} &= A_{\mu},
  \label{eq:floweq}\\
  G_{\mu\nu} &= \partial_{\mu}B_{\nu}-\partial_{\nu}B_{\mu}
  +[B_{\mu},B_{\nu}],
  &
  D_{\mu} &= \partial_{\mu}+[B_{\mu},\,\cdot\,],
\end{align}
where $A_{\mu}$ is the fundamental gauge field in QCD
(see Appendix~\ref{app:A} for unexplained notation; differentiation with
respect to the ``flow time'' $t$ is abbreviated by a dot).
Evidently, for increasing $t$ and as long as no singularities
develop, the flow equation~\eqref{eq:floweq} drives the gauge field
along the direction of steepest descent towards the
stationary points of the Yang--Mills action.

In lattice QCD, the simplest choice of the action
of the gauge field $U(x,\mu)$ is the Wilson action~\cite{Wilson}
\begin{equation}
  \Sw(U)=\frac{1}{g_0^2}\sum_p
  \Re\tr\{1-U(p)\},
  \label{eq:wact}
\end{equation}
where $g_0$ is the bare coupling,
$p$ runs over all oriented plaquettes on the lattice
and $U(p)$ denotes the product of
the link variables around $p$.
The associated flow $V_t(x,\mu)$ of
lattice gauge fields (the ``Wilson flow'')
is defined by the equations
\begin{equation}
  \dot{V}_t(x,\mu)=-g_0^2\left\{\partial_{x,\mu}\Sw(V_t)\right\}V_t(x,\mu),
  \qquad
  \left.V_t(x,\mu)\right|_{t=0}=U(x,\mu),
  \label{eq:wflow}
\end{equation}
in which $\partial_{x,\mu}$ stands for the natural $\suthree$-valued
differential operator with respect to the
link variable $V_t(x,\mu)$ (see Appendix~\ref{app:A}).
The existence, uniqueness and smoothness of the Wilson flow at
all positive and negative times $t$ is rigorously guaranteed on
a finite lattice~\cite{ArnoldI}. Moreover, from eq.~\eqref{eq:wflow}
one immediately concludes that
the action $\Sw(V_t)$ is a monotonically decreasing function of $t$.
The flow therefore tends to have a smoothing
effect on the field and it is, in fact, generated by infinitesimal
``stout link smearing'' steps~\cite{Stout}.

The Wilson flow previously appeared in ref.~\cite{TrivMaps}
in the context
of trivializing maps of field space.
Familiarity with this paper is not assumed, but some
mathematical results obtained there will be used here again.
An important goal in the following is to find out whether
the expectation values of local observables constructed
from the gauge field at positive flow time can be expected
to have a well-defined continuum limit.
Evidence for the existence of the limit is provided
by performing a sample calculation to one-loop order
of perturbation theory directly in the continuum theory, using
dimensional regularization, and through a
numerical study of the $\SUthree$ gauge theory at three values
of the lattice spacing.
Two applications of the Wilson flow are then discussed,
one concerning the scale-setting in lattice QCD and the
other the question of
how exactly the topological (instanton) sectors emerge
when the lattice spacing goes to zero.
